\noindent%
Comme annoncé en [\ref{1.29}], nous donnons deux calculs supplémentaires de
%
\begin{equation*}
  \Biggl(\Bigl(\bigl(1729^{1729} \bmod{17}\bigr)^{1729} \bmod{17}\Bigr)^{1729} \bmod{17}\Biggr) \cdots
\end{equation*}
%
en prouvant la congruence
%
\begin{equation*}
  12^{1729} \equiv 12 \pmod{17}.
\end{equation*}
%
Pour cela, nous exploitons la relation 
%
\begin{equation}
  1729 = 1^3 + 12^3.
\end{equation} 
%
% % % % % % % % % % % % % % % % % % % % % % % % % % % % % % % % % % % % % % % % % % % % % % % % % % % % % % % % % % % % % % % %
% Third variant: Freshman's dream
% % % % % % % % % % % % % % % % % % % % % % % % % % % % % % % % % % % % % % % % % % % % % % % % % % % % % % % % % % % % % % % %
\begin{proof}%
%\paragraph{} does not work. TODO:?
\textbf{À partir de l'égalité $1729 = 12^3 + 1^3$, par l'étude d'un polynôme unitaire de degré $3$.} 
\\
Commençons par remarquer que %
$12^3 + 1= 1729\equiv 12\pmod{17}$ implique 
%
\begin{equation}
  (12^3 + 1)^{17} \equiv 12^{17} \pmod{17}, 
\end{equation}
%
car \emph{$p$ est premier}.
Considérons alors le polynôme %
%
\begin{equation}
  P(X) \Def X^3 -X + 1 \in \Z[X].
\end{equation}
%
L'entier $12$ est une racine modulo $17$ de $P$. En effet, 
%
\begin{equation}
  P(12) = 12^3 - 12 + 1  \equiv 0 \pmod{17} 
\end{equation}
%
(ce qui établit au passage que $12^3 \equiv 11 \pmod{17}$). De surcroît, $17$ étant premier, 
%
\begin{equation}
  (12^3 + 1 )^{17} \equiv (12^3)^{17} + 1^{17} \pmod{17}.
\end{equation}
%
Combiné avec le fait que $(12^3 + 1)^{17} \equiv 12^{17}$, cela amène 
%
\begin{equation}
  (12^3)^{17} + 1  \equiv 12^{17} \pmod{17}.
\end{equation}
%
De façon équivalente, 
%
\begin{equation}
  P(12^{17}) = (12^{17})^3 - 12^{17} + 1 \equiv 0 \pmod{17}.
\end{equation}
%
L'entier $12^{17}$ est donc également une racine modulo $17$ de $P$. Nous montrons que c'est la seule, et, par voie de %
conséquence, que 
%
\begin{equation}
  12^{17} \equiv 12 \pmod{17}. 
\end{equation}
%
Comme 
%
\begin{align}
  P(n) &= n^3 - n  + 1 \\
  &= n \cdot (n^2 - 1) + 1 \\
  &= n \cdot (n-1) \cdot (n+1) +1  
\end{align}
%
pour tout entier $n$, résoudre la congruence 
%
\begin{equation}
  P(n) \equiv 0 \pmod{17} 
\end{equation}
%
revient à résoudre 
%
\begin{equation}
  (n-1) \cdot n \cdot (n+1) \equiv -1 \pmod{17}.
\end{equation}
%
Seuls les $n \equiv 12 \pmod{17}$ satisfont cette relation, ce que montre une vérification excluant les cas triviaux
%
\begin{align}
  \text{\texttt{(1*2*3 + 1) \% 17 }}&\text{\texttt{== 0    \# False}}\\
  &\hspace{0.14em}\vdots \nonumber\\
  \text{\texttt{(14*15*16 + 1) \% 17 }}&\text{\texttt{== 0    \# False.}}
\end{align}
%
L'entier $12$ étant l'unique racine modulo $17$ de $P$, nous avons bien établi que $12^{17} \equiv 12 \pmod{17}$. Utilisant %
cette congruence avec 
%
\begin{align}
  1729 &= 1717 + 12 \\
  &= 17 \times 101  + 12 \\
  &= 17 \times (5 \times 17 + 16) + 12 \\
  &= 17^2 \times 5 + 17 \times 16 + 12, 
\end{align}
%
nous obtenons 
%
\begin{align}
  12^{1729} &= ((12^{17})^{17})^5 \times (12^{17})^{16} \times 12^{12} \\
  &\equiv 12^5 \times 12^{16} \times 12^{12} \\
  &\equiv 12^{5+12} \times 12^{16} \\
  &\equiv 12^1 \times 12^{16} \\
  &\equiv 12 \pmod{17},
\end{align}
%
ce qui termine cette troisième démonstration.
%
% % % % % % % % % % % % % % % % % % % % % % % % % % % % % % % % % % % % % % % % % % % % % % % % % % % % % % % % % % % % % % % %
% Fourth variant: Freshman's dream, with the delta
% % % % % % % % % % % % % % % % % % % % % % % % % % % % % % % % % % % % % % % % % % % % % % % % % % % % % % % % % % % % % % % %
\paragraph{Variante: toujours à partir de l'égalité $1729 = 12^3 + 1^3$ mais avec la recherche d'un résidu quadratique} %
Cette preuve ne diffère de la précédente que dans la manière d'établir que $12$ est la seule racine modulo $17$ de $P$. Ici, %
nous cherchons un polynôme 
%
\begin{equation}
  Q(X)=X^2+ b X+ c \in \Z[X]
\end{equation}
%
vérifiant
%
\begin{equation}
  P(n) \equiv  (n-12)\cdot Q(n) \pmod{17} \qquad (n \in \Z).
\end{equation}
%
Le développement $P(n) \equiv  n^3 + b n^2 + c n - 12 n^2 - 12 bn - 12 c \pmod{17}$ donne %
$b\equiv 12 \pmod{17}$ et $c \equiv 7 \pmod{17}$. Déterminons l'ensemble des racines modulo $17$ de $Q$, c'est-à-dire tout %
entier relatif $n$ vérifiant%
%
\begin{align}
  Q(n) &\Def  n^2 + 2\cdot 6 \cdot n + 7 \\
  &= (n + 6)^2 - 29 \\
  &\equiv (n + 6)^2 - 12 \pmod{17}.
\end{align}
%
De façon équivalente, il s'agit de résoudre la congruence
%
\begin{equation}
  Q(n) \equiv (n + 6)^2 - r^2 \qquad (r^2 \equiv 12 \pmod{17})
\end{equation}
%
avec $0\leq r < 17$. À cette fin nous calculons modulo $17$ les résidus $k^2$\;\,($0 \leq k < 17)$ et constatons que $12$ %
n'en est pas un (les résidus calculés sont $0, 1, 2, 4, 8, 9, 13, 15, 16$). Ainsi, $12$ n'est pas un \emph{résidu quadratique} 
modulo $17$, et $Q$ n'a donc pas de racine modulo $17$. En conclusion, l'entier $12$ est la seule racine modulo $17$ de $P$, %
ce qui achève la démonstration.
\end{proof}